\section{Related Work}
\label{sec:related}

\paragraph{OS VM compression}
Operating systems compress inactive pages as a middle ground between
keeping them resident and swapping to disk.  macOS compresses inactive
pages using WKdm~\cite{wilson1999case,apple2013compressed}, a
word-granularity dictionary compressor designed for 4\,KiB
pointer-heavy pages\punt{; on 16\,KiB ARM64 pages with text or serialized
data, WKdm achieves poor ratios (98.9\% on JSON, 93.7\% on key-value
records; Table~\ref{tab:algo_compare})}.  Linux
zswap~\cite{jennings2013zswap} intercepts pages heading to swap and
compresses them with LZ4, LZO, or zstd, operating per-page with no
cross-page context and no knowledge of which bytes are live objects;
zram~\cite{zram2024} provides a compressed block device in RAM\@.
Windows~10 compresses pages from the modified page list using the
Xpress algorithm~\cite{microsoft2015compression}.  All three
mechanisms treat each page as an opaque byte array.  \sys operates
above the OS, exploiting object liveness, call-site patterns,
and per-size-class policies for higher~ratios.

\paragraph{Userspace memory compression}
Chen et al.~\cite{chen2003heap} compress the Java heap during garbage
collection.  Their Mark-Compact-Compress (MCC) collector triggers
compression reactively when allocation fails and compaction alone
cannot free enough space; it then compresses \emph{all} live objects
using per-object zero-byte removal (a bitmap records which bytes are
non-zero).  Objects are accessed through a handle table; decompressing
a compressed object requires allocating a fresh block, restoring the
data, and updating the handle.  By contrast, \sys compresses
\emph{proactively} on cold pages via a concurrent background thread,
uses general-purpose compressors (LZ4/zstd) that exploit repeated
patterns beyond just zero bytes, and decompresses transparently through
virtual-memory faults with no handle indirection on the normal access
path.
Rumble~\cite{rumble2014log} describes a log-structured allocator for
RAMCloud that compresses objects inline.  Unlike \sys, Rumble requires
application changes to use its API and does not support transparent
\texttt{malloc}/\texttt{free} interposition.
Tsai and Sanchez~\cite{tsai2019compress} propose Zippads, an
object-granularity compressed memory hierarchy that compresses objects
rather than cache lines, achieving higher compression ratios by
exploiting object structure.  Zippads requires hardware support
(a compressed object cache and modified TLBs); \sys achieves similar
object-aware benefits in software by exploiting allocator metadata.
Plumber~\cite{ayers2019asmdb} focuses on detecting and eliminating
memory waste rather than compressing cold data.

\paragraph{Memory-efficient allocators}
Mesh~\cite{powers2019mesh} reduces fragmentation by physically
merging pages with non-overlapping live objects via virtual memory
remapping.  In our evaluation (\S\ref{sec:eval:apps}), Mesh
reduces RSS on workloads with sparse pages (35\% on RocksDB, 42\%
on DuckDB) but not on workloads where pages remain fully occupied.
The two techniques are complementary: Mesh eliminates internal
waste, while \sys compresses cold data.
Hoard~\cite{berger2000hoard} and
jemalloc~\cite{evans2015jemalloc} use arenas for thread
scalability; \sys repurposes arenas for data homogeneity.
tcmalloc~\cite{google2024tcmalloc} uses per-CPU caches and size
classes similar to \sys but does not compress cold pages.
mimalloc~\cite{leijen2019mimalloc} achieves excellent performance
through segment-based allocation with free list sharding.

\paragraph{Tiered and disaggregated memory}
CXL-attached memory~\cite{pond2022cxl, maruf2023tpp} and
software-defined far memory~\cite{amaro2023memory} extend available
memory at higher access latency.  TMO~\cite{weiner2022tmo}
transparently offloads cold pages to compressed swap or secondary
storage at Meta's scale, using PSI (pressure stall information) to
balance memory savings against application performance.  TMO operates
at the OS level on whole pages; \sys operates within the allocator,
exploiting object liveness and structure for higher ratios.  The two
are complementary: \sys could compress locally before TMO offloads
remaining cold~pages.

\paragraph{Compressed caches and memory}
LCP~\cite{pekhimenko2012linearly} and BDI~\cite{pekhimenko2012linearly}
compress cache lines at the hardware level.  Linearly compressed
pages~\cite{pekhimenko2012linearly} pack compressed cache lines
into fewer physical pages.  MemSys~\cite{vijaykumar2015case}
explores compressed main memory with hardware support.
These hardware approaches target
last-level cache or DRAM and require architecture changes, while \sys
operates in userspace software.

\paragraph{Pointer compression and managed heaps}
V8's Oilpan garbage collector~\cite{nickel2024oilpan} uses pointer
compression to reduce pointer sizes from 64 to 32 bits within a 4\,GiB
heap cage, saving $\sim$16\% of managed-heap memory.  This reduces
memory used by pointers but does not compress data content.
Apache Ignite~\cite{ignite2024compression} compresses data pages when
flushing to persistent storage using LZ4 or zstd, achieving 2--4$\times$
compression on typical workloads; however, pages remain uncompressed in
memory.  \sys compresses pages \emph{in memory}, reducing RSS without
relying on disk-backed persistence.

\paragraph{Zero-page optimization}
The Linux kernel's KSM (Kernel Same-page Merging)~\cite{arcangeli2009ksm}
deduplicates identical pages, including zero pages.
Chow et al.~\cite{chow2005shredding} proposed \emph{secure
deallocation}---zeroing freed memory to limit the lifetime of
sensitive data.  \sys's zero-on-free serves a dual purpose:
creating compressible zero runs within partially-filled pages
while also limiting data lifetime.
